\documentclass[11pt,a4paper]{article}
\usepackage[margin=1in]{geometry}
\usepackage[T1]{fontenc}
\usepackage{lmodern,microtype}
\usepackage{amsmath,amssymb,amsthm,mathtools,mathrsfs}
\usepackage{booktabs,array,longtable,tabularx}
\usepackage{enumitem}
\usepackage{xcolor}
\usepackage{hyperref}
\hypersetup{colorlinks=true,linkcolor=blue!45!black,citecolor=blue!45!black,
  urlcolor=blue!45!black,
  pdftitle={Optimal Fresh-Oracle Complexity for Nonconvex-PL Minimax Optimization},
  pdfauthor={Colombo Team},
  pdfkeywords={minimax optimization, Polyak-Lojasiewicz condition, stochastic oracle, randomized lower bound}}
\setlength{\parindent}{1.25em}
\setlength{\parskip}{0.15em}
\emergencystretch=2em
\allowdisplaybreaks[1]
\setlist[itemize]{leftmargin=*,itemsep=0.3em,topsep=0.4em}
\setlist[enumerate]{leftmargin=*,itemsep=0.3em,topsep=0.4em}
\newcommand{\R}{\mathbb R}
\newcommand{\E}{\mathbb E}
\newcommand{\Pp}{\mathbb P}
\newcommand{\eps}{\varepsilon}
\newcommand{\norm}[1]{\left\lVert #1\right\rVert}
\newcommand{\ip}[2]{\left\langle #1,#2\right\rangle}
\newcommand{\one}{\mathbf 1}
\newcommand{\ind}{\mathbf 1}
\newcommand{\cA}{\mathcal A}
\newcommand{\cF}{\mathcal F}
\newcommand{\cR}{\mathcal R}
\newcommand{\rprog}{\operatorname{rprog}}
\newcommand{\prog}{\operatorname{prog}}
\newcommand{\diag}{\operatorname{diag}}
\newcommand{\Lip}{\operatorname{Lip}}
\newcommand{\supp}{\operatorname{supp}}
\newcommand{\res}{\mathcal S_\kappa}
\DeclareMathOperator*{\argmax}{arg\,max}
\DeclareMathOperator*{\argmin}{arg\,min}
\newtheorem{theorem}{Theorem}[section]
\newtheorem{lemma}[theorem]{Lemma}
\newtheorem{proposition}[theorem]{Proposition}
\newtheorem{corollary}[theorem]{Corollary}
\theoremstyle{definition}
\newtheorem{definition}[theorem]{Definition}
\theoremstyle{remark}
\newtheorem{remark}[theorem]{Remark}
\numberwithin{equation}{section}
\title{Optimal Fresh-Oracle Complexity for\\Nonconvex--P\L{} Minimax Optimization}
\author{Colombo Team\\[3pt]
\normalsize\href{mailto:colombo.2026@outlook.com}{\texttt{colombo.2026@outlook.com}}}
\date{}
\begin{document}
\maketitle

\begin{abstract}
We study the stochastic first-order oracle complexity of smooth nonconvex
minimax optimization under a Polyak--\L{}ojasiewicz condition in the
maximization variable. The oracle returns an unbiased joint-gradient
estimate with variance at most $\sigma^2$, using an independent fresh seed
at each call. For the weighted stationarity criterion
$\E[\norm{\nabla_x f}^2+\kappa\norm{\nabla_y f}^2]\le\eps^2$,
where $\kappa=L/\mu$, we prove the lower bound
\[
 \Omega\!\left(\frac{L\kappa H}{\eps^2}
       +\frac{L\kappa^2H\sigma^2}{\eps^4}\right).
\]
Here $H$ bounds the initial value gap together with an auxiliary
primal--dual gap. The result holds for arbitrary adaptive randomized
algorithms, with unrestricted queries and outputs, when $\kappa$ is
larger than a universal constant and $\eps^2\le cLH$.
It matches the stochastic Smoothed-AGDA upper bound of Yang et al.
under the same stationarity criterion and initialization budget.
The proof extends the deterministic weighted saddle chain of Pan and Li
using a globally smooth cap that preserves the dual P\L{} inequality,
Bernoulli masking of complete gradient increments, and independently
rotated hidden frames. Prefix-truncated functions give an ideal oracle
whose responses depend only on revealed directions; an exact coupling
then transfers the information bound to a fixed oracle on a fixed
objective. The construction also provides explicit smoothness,
P\L{}, scaling, and dimension bounds.
\end{abstract}

\noindent\textbf{Keywords.}
Minimax optimization; Polyak--\L{}ojasiewicz condition; stochastic first-order
oracles; randomized lower bounds; oracle complexity.

\begin{samepage}
\noindent\textbf{AI Usage.}
Nearly the entire research pipeline for this paper was carried out by
Colombo, powered by GPT-5.6 Sol.
\par
\end{samepage}

\section{Introduction}
\label{sec:introduction}

Consider the minimax problem
\[
 \min_{x\in\R^{d_x}}\Phi(x),\qquad
 \Phi(x):=\max_{y\in\R^{d_y}}f(x,y),
\]
where $f$ has an $L$-Lipschitz gradient and its dependence on $y$ satisfies
a $\mu$-Polyak--\L{}ojasiewicz (P\L{}) inequality. The outer objective may
be nonconvex, and the inner problem need not be concave. The P\L{}
condition nevertheless relates the inner objective gap to the dual
gradient, making quantitative stationarity guarantees possible.
The resulting nonconvex--P\L{} (NC--P\L{}) class raises a basic complexity
question: how much of the dependence on the condition parameter
$\kappa:=L/\mu$ is unavoidable when an algorithm receives noisy gradients?

We answer this question for the weighted joint-gradient residual
\begin{equation}
 \res(x,y):=\norm{\nabla_x f(x,y)}^2
                +\kappa\norm{\nabla_y f(x,y)}^2.
 \label{eq:residual}
\end{equation}
The target is $\E\res(\widehat x,\widehat y)\le\eps^2$.
We charge both the initial value gap and the auxiliary gap used in the
analysis of stochastic Smoothed-AGDA to a single budget $H$; the precise
definition is given in Section~\ref{sec:model}. In a fresh oracle model,
each call returns one unbiased estimate of the joint gradient, and the
random seed is independent of all previous calls.

\paragraph{Main result.}
For sufficiently large $\kappa$ and sufficiently small
$\eps^2/(LH)$, both in the sense of universal numerical thresholds,
the minimax number of fresh oracle calls is
\begin{equation}
 \Theta\!\left(
 \frac{L\kappa H}{\eps^2}
 +\frac{L\kappa^2H\sigma^2}{\eps^4}
 \right).
 \label{eq:intro-rate}
\end{equation}
The lower bound allows arbitrary randomized algorithms, including
methods that query outside the span of earlier responses or return an
unqueried point. The hard objective and oracle are fixed before the run.
The dimensions are finite and may grow with the query budget. Thus the
claim concerns the worst case over dimensions; it does not assert an
optimal rate in a prescribed fixed dimension.

\paragraph{Why randomized lifting requires additional structure.}
A deterministic chain can hide the next useful coordinate from a
method that respects the coordinates already exposed. An unrestricted
algorithm can instead probe arbitrary directions, and stochastic
responses must remain unbiased at every query. A random rotation alone
does not resolve either issue: sufficiently large queries can expose
hidden directions, and masking a single dual component can leave the
same direction encoded in another component of the full gradient.
The construction below addresses both difficulties while preserving
global smoothness and the inner P\L{} inequality.

The proof uses three ingredients. First, a radial cap makes each dual
block independent of its hidden frame outside a bounded cone, while
retaining its exact maximum and a uniform P\L{} estimate. Second, the
oracle masks the gradient of the entire scalar increment associated
with the next coordinate. Its norm is $O(\eps/\sqrt\kappa)$ after
scaling, which produces the extra factor of $\kappa$ in the stochastic
term. Third, each ideal response is defined by a prefix-truncated
function before any concentration event is imposed. This makes the
conditional distribution of unrevealed directions explicit and permits
an exact coupling to the fixed oracle.

\subsection{Relation to previous work}
\label{sec:related}

Yang et al.~\cite{yang} established the stochastic Smoothed-AGDA upper
bound for the weighted residual in \eqref{eq:residual}. We use their
potential estimate with the complete initialization budget $H$.
Pan and Li~\cite{panli} proved a deterministic lower bound with linear
dependence on $\kappa$ for value-gradient stationarity, under a bound on
the initial value gap alone. Table~\ref{tab:comparison} records the
different targets and budgets explicitly.

\begin{table}[t]
\centering
\caption{Related complexity bounds for smooth NC--P\L{} problems.
In the first and third rows the target is $\E\res\le\eps^2$ and the
budget is the charged initialization $H$ in \eqref{eq:charged-init};
these two rows match. In the middle row,
the budget is $\Phi(0)-\inf\Phi\le\Delta$ and the target is the value gradient.
All lower bounds are worst-case bounds over finite dimensions.}
\label{tab:comparison}
\small
\renewcommand{\arraystretch}{1.18}
\begin{tabularx}{\textwidth}{@{}>{\raggedright\arraybackslash}p{0.19\textwidth}
 >{\raggedright\arraybackslash}p{0.22\textwidth}
 >{\raggedright\arraybackslash}p{0.20\textwidth}
 >{\raggedright\arraybackslash}X@{}}
\toprule
Work and guarantee & Oracle and algorithm & Target and budget & Complexity and regime\\
\midrule
Yang et al.~\cite{yang}, upper bound & Fresh stochastic gradients;
Smoothed-AGDA & $\E\res\le\eps^2$; charged $H$
 & $O\!\left(1+\dfrac{L\kappa H}{\eps^2}
 +\dfrac{L\kappa^2H\sigma^2}{\eps^4}\right)$\par
 $\kappa\ge1$, $\eps>0$ \\
\addlinespace[4pt]
Pan--Li~\cite{panli}, lower bound & Exact first-order oracle;
arbitrary deterministic methods & $\norm{\nabla\Phi}\le\eps$; value gap $\Delta$
 & $\Omega\!\left(\dfrac{L\kappa\Delta}{\eps^2}\right)$\par
 $\kappa\ge c_0$, $\eps^2\le c_1L\Delta$ \\
\addlinespace[4pt]
This work, lower bound & Fresh stochastic gradients;
arbitrary randomized methods & $\E\res\le\eps^2$; charged $H$
 & $\Omega\!\left(\dfrac{L\kappa H}{\eps^2}
 +\dfrac{L\kappa^2H\sigma^2}{\eps^4}\right)$\par
 $\kappa\ge\kappa_0$, $\eps^2\le c_\eps LH$ \\
\bottomrule
\end{tabularx}
\end{table}

Our deterministic starting point is the weighted dual chain and
perspective coupling of Pan and Li~\cite{panli}. We use the outer chain
and randomized-embedding framework of Carmon et al.~\cite{carmon},
together with the normalization and soft
projection recorded by Arjevani et al.~\cite{arjevani}, whose
probabilistic zero-chain construction also motivates the Bernoulli
masking. The additional work here is to make these components compatible
with a globally valid fresh oracle and an unrestricted randomized
algorithm, while proving the weighted joint-gradient obstruction.
The cap, complete-increment estimate, and prefix-based coupling are
proved explicitly below.

The stationarity distinction matters. A guarantee for
$\res(x,y)$ is a statement about the two gradients of $f$ at one pair;
a guarantee for $\norm{\nabla\Phi(x)}$ concerns the maximized value
function. Our lower-bound construction has an explicit comparison
between these quantities, but that comparison is not assumed for the
whole function class. The matching conclusion in
\eqref{eq:intro-rate} therefore concerns the weighted joint residual.

\paragraph{Organization.}
Section~\ref{sec:model} specifies the model and states the matching
bounds. Section~\ref{sec:overview} explains the parameter mechanism.
Section~\ref{sec:construction} constructs the globally smooth saddle
family. Sections~\ref{sec:prefix} and~\ref{sec:scaling} establish its
information structure and scaling. Sections~\ref{sec:oracle}
and~\ref{sec:haar} define the oracle and complete the randomized lower
bound. The appendices collect auxiliary derivative estimates and
explicit constants.

\section{Oracle model and matching bounds}
\label{sec:model}

All norms are Euclidean; matrix norms are the corresponding operator
norms. A continuously differentiable function is $L$-smooth if its
gradient is globally $L$-Lipschitz. We write $C^{1,1}$ for this regularity
when no particular Lipschitz constant is specified. Constants in
$O(\cdot)$, $\Omega(\cdot)$, and $\Theta(\cdot)$ are numerical and
independent of all problem parameters.

\subsection{Function class, oracle, and initialization}

Let $L,H>0$, $\sigma\ge0$, and $\kappa\ge1$, and set $\mu=L/\kappa$.
The objective $f:\R^{d_x}\times\R^{d_y}\to\R$ is jointly $L$-smooth.
For every $x$, the maximum
$\Phi(x):=\max_y f(x,y)$ is finite and attained. We assume
$\inf_x\Phi(x)>-\infty$ and the dual P\L{} inequality
\begin{equation}
 \norm{\nabla_y f(x,y)}^2
 \ge 2\mu\bigl(\Phi(x)-f(x,y)\bigr)
 \qquad\text{for all }(x,y).
 \label{eq:model-pl}
\end{equation}
The parameters $L$ and $\mu$ are admissible bounds; they need not be
the smallest smoothness constant or the largest P\L{} modulus.

A fresh stochastic first-order oracle is a fixed measurable map
$G_f(z;\xi)$, where $z=(x,y)$, satisfying
\begin{equation}
 \E_\xi G_f(z;\xi)=\nabla f(z),\qquad
 \E_\xi\norm{G_f(z;\xi)-\nabla f(z)}^2\le\sigma^2
 \quad\text{for every }z.
 \label{eq:model-oracle}
\end{equation}
Different calls use independent seeds with the same fixed law, also
independent of the algorithm's internal randomness. One call returns
both gradient blocks at one point. Obtaining the blocks at two different
points costs two calls. The oracle supplies no function values, and
seeds cannot be reused at a second point. No smoothness or
conservativity of the individual random vector field $G_f(\cdot;\xi)$
is required. In particular, this model does not assume that each
response is the gradient of a globally defined sample loss.

The algorithm may choose arbitrary adaptive query points, perform
unrestricted computation between calls, and return any measurable
output $(\widehat x,\widehat y)$, including a point that it has never
queried. The output requirement is
\begin{equation}
 \E\res(\widehat x,\widehat y)
 =\E\!\left[
 \norm{\nabla_x f(\widehat x,\widehat y)}^2
 +\kappa\norm{\nabla_y f(\widehat x,\widehat y)}^2
 \right]\le\eps^2,
 \label{eq:output-criterion}
\end{equation}
where the expectation includes both sources of randomness.

Following the potential used by Yang et al.~\cite{yang}, we supply
initial points $x_0,z_0\in\R^{d_x}$ and $y_0\in\R^{d_y}$ and impose
\begin{equation}
 \begin{gathered}
 \Phi(z_0)-\inf\Phi+b_0\le H,\\
 b_0=2\left[
   \sup_y\widehat f_{z_0}(x_0,y)
   -\inf_x\widehat f_{z_0}(x,y_0)
 \right],\qquad
 \widehat f_{z_0}(x,y)=f(x,y)+L\norm{x-z_0}^2.
 \end{gathered}
 \label{eq:charged-init}
\end{equation}
We call this the \emph{charged initialization budget}. It includes the
cost represented by the initial auxiliary primal--dual gap, rather than
bounding only the value gap. The lower-bound instances use
$x_0=z_0=0$ and $y_0=0$; algorithms remain free to make their first
query elsewhere.

Let $\mathsf Q^*(\eps;L,\kappa,H,\sigma)$ be the least integer $Q$
for which a randomized algorithm, operating in every finite pair of
dimensions, uses at most $Q$ calls on every run and guarantees
\eqref{eq:output-criterion} for every admissible objective and oracle.
The query bound is a fixed horizon, not a bound only on the expected
stopping time.

\subsection{Main lower bound}

\begin{theorem}[Randomized lower bound]
\label{thm:main-lower}
There are universal constants $c,c_\eps>0$ and $\kappa_0\ge1$ with
the following property. Let $L,H>0$, $\sigma\ge0$,
$\kappa\ge\kappa_0$, and $0<\eps^2\le c_\eps LH$.
For every adaptive randomized algorithm using at most
\[
 Q\le c\left(
   \frac{L\kappa H}{\eps^2}
   +\frac{L\kappa^2H\sigma^2}{\eps^4}
 \right)
\]
fresh oracle calls, there exist finite dimensions, an admissible,
jointly $L$-smooth deterministic objective satisfying
\eqref{eq:model-pl} and \eqref{eq:charged-init},
and a fixed oracle satisfying \eqref{eq:model-oracle}, such that
\[
 \E\res(\widehat x,\widehat y)>\eps^2.
\]
The objective and oracle may depend on the algorithm and its query
horizon, but neither depends on the realized transcript or the realized
internal random seed.
\end{theorem}

The theorem is proved in Section~\ref{sec:completion}. Explicit valid
choices of $\kappa_0$ and $c_\eps$ are given in
Section~\ref{sec:scaling}. Their numerical values are conservative;
the result concerns the parameter exponents and matching absolute-factor
complexity in the stated regime.

\subsection{Upper bound under the same budget}

For completeness, we spell out how the known upper bound uses
\eqref{eq:charged-init}. Stochastic Smoothed-AGDA~\cite{yang} maintains
an anchor $z_t$ and performs
\begin{equation}
 \begin{aligned}
 x_{t+1}&=x_t-\eta\bigl(G_x(x_t,y_t;\xi_t^x)+2L(x_t-z_t)\bigr),\\
 y_{t+1}&=y_t+\frac{\eta}{48}G_y(x_{t+1},y_t;\xi_t^y),\\
 z_{t+1}&=z_t+\beta(x_{t+1}-z_t),
 \qquad \beta=\frac{\eta\mu}{48\cdot1600}.
 \end{aligned}
 \label{eq:upper-algorithm}
\end{equation}
The two seeds at iteration $t$ are independent and fresh.

\begin{proposition}[Smoothed-AGDA upper bound]
\label{prop:upper-bound}
There is a universal constant $C\ge1$ such that, for every
$0<\eta\le1/(3L)$, the iterates in \eqref{eq:upper-algorithm} satisfy
\begin{equation}
 \frac1T\sum_{t=0}^{T-1}\E\res(x_t,y_t)
 \le C\kappa\left(\frac{H}{\eta T}+L\eta\sigma^2\right).
 \label{eq:upper-potential}
\end{equation}
Consequently, a uniformly random iterate satisfies
\eqref{eq:output-criterion} within
\[
 O\!\left(1+\frac{L\kappa H}{\eps^2}
              +\frac{L\kappa^2H\sigma^2}{\eps^4}\right)
\]
fresh oracle calls.
\end{proposition}

\begin{proof}
The potential argument in the proof of Theorem~4.1 of
Yang et al.~\cite{yang}, Supplement~C, immediately after its
equation~(41), gives
\eqref{eq:upper-potential}: its initial potential gap is at most
$\Phi(z_0)-\inf\Phi+b_0\le H$.
Our joint smoothness and variance assumptions imply the corresponding
blockwise assumptions in that argument. Choose
\[
 \eta=\min\left\{\frac1{3L},
       \frac{\eps^2}{2C\kappa L\sigma^2}\right\},\qquad
 T=\left\lceil\frac{2C\kappa H}{\eta\eps^2}\right\rceil,
\]
omitting the second entry of the minimum when $\sigma=0$.
The two terms in \eqref{eq:upper-potential} are at most
$\eps^2/2$ each. A uniform output index converts their average into
the output expectation, and every iteration costs two calls.
\end{proof}

\begin{corollary}[Matching fresh-oracle complexity]
\label{cor:matching}
Under the parameter conditions of Theorem~\ref{thm:main-lower},
\begin{equation}
 \mathsf Q^*(\eps;L,\kappa,H,\sigma)
 =\Theta\!\left(\frac{L\kappa H}{\eps^2}
                  +\frac{L\kappa^2H\sigma^2}{\eps^4}\right).
 \label{eq:matching-rate}
\end{equation}
\end{corollary}

\begin{proof}
Combine Theorem~\ref{thm:main-lower} and
Proposition~\ref{prop:upper-bound}. The additive constant in the upper
bound is absorbed by the first term in the stated regime.
\end{proof}

The ratio of the stochastic term to the deterministic term is
$\kappa\sigma^2/\eps^2$. Thus the deterministic rate governs when
$\sigma^2\lesssim\eps^2/\kappa$, and the stochastic term governs
when $\sigma^2\gtrsim\eps^2/\kappa$. The theorem covers both regimes,
including the exact-gradient case $\sigma=0$.

\section{Proof strategy and the source of the rate}
\label{sec:overview}

The construction has $M$ primal links, each preceded by a dual chain
of length $N$. The lengths will satisfy
\[
 M\asymp\frac{LH}{\eps^2},\qquad N\asymp\kappa,
\]
with small universal proportionality constants. Maximizing each dual
block recovers one link of the Carmon chain exactly. The initial gap
is therefore $O(M)$ before scaling, whereas the dual P\L{} modulus
is of order $1/N$. A local-deficit estimate transfers the outer
gradient obstruction to the weighted residual in
\eqref{eq:residual}.

The information order is
\[
 \underbrace{y^{(1)}_1,\ldots,y^{(1)}_N}_{\text{first dual block}},x_1,
 \underbrace{y^{(2)}_1,\ldots,y^{(2)}_N}_{\text{second dual block}},x_2,
 \ldots,x_M.
\]
These are latent coordinates in independently rotated frames. The
terminal dual coordinate opens the next primal link, and that primal
link activates the next dual block. Flat regions in the scalar gates
ensure that one response cannot make a two-coordinate jump at either
boundary.

The first $N-1$ coordinates of each dual block have a complete
gradient increment of norm at most
$\Gamma=64\eps/\sqrt N$. The fixed oracle reveals that increment
with probability
\[
 p=\frac{\Gamma^2}{\Gamma^2+\sigma^2},
 \qquad
 \frac1p=1+\frac{N\sigma^2}{64^2\eps^2}.
\]
The resulting variance is at most $\sigma^2$ everywhere. Reaching the
last primal direction requires $M(N-1)$ such revelations, giving
\begin{equation}
 \frac{M(N-1)}p
 \asymp MN+\frac{MN^2\sigma^2}{\eps^2}
 \asymp\frac{L\kappa H}{\eps^2}
          +\frac{L\kappa^2H\sigma^2}{\eps^4}.
 \label{eq:rate-mechanism}
\end{equation}
This calculation identifies the two distinct sources of condition
dependence: $N$ coordinates must be traversed at each primal link,
and the variance constraint introduces a further factor $N$ in the
waiting time for a revelation.

Unrestricted queries are handled on both sides of the saddle.
A soft projection bounds the effective primal query, while the dual
cap removes all hidden-frame dependence outside a bounded cone.
Random-frame concentration then controls the remaining queries.
The proof includes the algorithm's final output as one additional
projection test. Prefix-truncated functions make this concentration
argument valid for an adaptive transcript without allowing the true
oracle to depend on that transcript.

\section{A globally smooth saddle family}
\label{sec:construction}

This section constructs an unscaled saddle function with $M$ primal
links and $M$ dual blocks of length $N$. The primal ambient dimension
is $d_x$; each dual block has ambient dimension $d_b$, so the total
dual dimension is $Md_b$. The hidden frames are
$U\in\R^{d_x\times M}$ and
$V_i\in\R^{d_b\times N}$, all with orthonormal columns.
They are fixed throughout this section and randomized only in
Section~\ref{sec:haar}.

\subsection{Outer geometry and soft projection}
\label{sec:outer}

We use the Carmon chain~\cite{carmon} in the normalization of
Arjevani et al.~\cite{arjevani}. Define
\[
 \psi(t)=
 \begin{cases}
 0,&t\le\frac12,\\
 \exp\!\left(1-\dfrac1{(2t-1)^2}\right),&t>\frac12,
 \end{cases}
 \qquad
 \phi(t)=\sqrt e\int_{-\infty}^t e^{-u^2/2}\,du,
\]
and
\begin{equation}
 h(s,t)=\psi(-s)\phi(-t)-\psi(s)\phi(t),\qquad
 F_M(s)=\sum_{i=1}^M h(s_{i-1},s_i),\quad s_0=1.
 \label{eq:outer-chain}
\end{equation}
For a vector $s\in\R^M$ and a threshold $a\ge0$, write
\[
 \prog_a(s)=\max\bigl(\{i\in\{1,\ldots,M\}:|s_i|>a\}\cup\{0\}\bigr).
\]
The standard chain estimates are
\begin{equation}
 F_M(0)-\inf F_M\le12M,\qquad
 \Lip(\nabla F_M)\le152,\qquad
 \norm{\nabla F_M(s)}_\infty\le23.
 \label{eq:carmon-basic}
\end{equation}
These numerical bounds are recorded in Lemma~2 of
Arjevani et al.~\cite{arjevani}. We also use a direct sign observation:
every first partial derivative of $h$ is nonpositive. If $j$ is the
first index with $|s_j|<1$, then $|s_{j-1}|\ge1$ and
\begin{equation}
 [\nabla F_M(s)]_j\le-\phi'(s_j)<-1.
 \label{eq:outer-sign}
\end{equation}
Indeed, the incoming term has magnitude at least $\phi'(s_j)$, and
the outgoing term has the same sign.

For $R_x=230\sqrt M$, define the soft projection and regularized value
function
\begin{equation}
 \varrho(x)=\frac{x}{\sqrt{1+\norm{x}^2/R_x^2}},\qquad
 \widehat\Phi_{M,U}(x)=F_M(U^\top\varrho(x))+
                           \frac1{10}\norm x^2.
 \label{eq:outer-soft}
\end{equation}
This is the soft-projection construction of
Arjevani et al.~\cite{arjevani}, equation~(27). We give the obstruction
argument to make its numerical normalization explicit.

\begin{lemma}[Gradient obstruction after soft projection]
\label{lem:soft-projection}
For every $x$,
\begin{equation}
 \prog_{1/4}(U^\top\varrho(x))<M
 \quad\Longrightarrow\quad
 \norm{\nabla\widehat\Phi_{M,U}(x)}\ge\frac12.
 \label{eq:soft-obstruction}
\end{equation}
\end{lemma}

\begin{proof}
Put $R=R_x$, $r=\norm x$,
$\alpha=(1+r^2/R^2)^{-1/2}$, $s=U^\top\alpha x$, and
$g=\nabla F_M(s)$. Differentiation gives
\[
 \nabla\widehat\Phi(x)
 =\alpha Ug+
 \left(\frac15-\alpha\frac{\ip{x}{Ug}}{R^2+r^2}\right)x.
\]
The Jacobian of $\varrho$ has norm at most $\alpha$.
If $r\ge R/2$, then
\[
 \norm{\nabla\widehat\Phi(x)}
 \ge\frac r5-\alpha\,23\sqrt M
 \ge23\sqrt M\left(1-\frac2{\sqrt5}\right)>\frac12.
\]
For $r\le R/2$, we have $\alpha\ge2/\sqrt5$ and
\[
 c:=\frac{\ip{x}{Ug}}{R^2+r^2},\qquad
 |c|\le\frac{\norm g}{2R}\le\frac1{20}.
\]
Thus $b:=1/(5\alpha)-c$ is positive and at most
$\sqrt5/10+1/20$. The hypothesis gives $|s_M|\le1/4$,
so there is a first $j$ with $|s_j|<1$. By
\eqref{eq:outer-sign},
\[
 \ip{u_j}{\nabla\widehat\Phi(x)}
 =\alpha g_j+bs_j
 <-\frac2{\sqrt5}+\frac{\sqrt5}{10}+\frac1{20}
 <-\frac12.
\]
This proves the claim in both cases.
\end{proof}

To express each link as a saddle block, use the scale and nonnegative
target from Pan and Li~\cite{panli}:
\begin{equation}
 \rho(s)=\sqrt{\psi(s)+\psi(-s)},\qquad
 \widetilde h(s,t)=h(s,t)+5\rho(s)^2.
 \label{eq:outer-rho-target}
\end{equation}
They satisfy
\begin{equation}
 \widetilde h(s,t)=
 \begin{cases}
  \rho(s)^2(5-\phi(t)),&s>1/2,\\
  0,&|s|\le1/2,\\
  \rho(s)^2(5+\phi(-t)),&s<-1/2.
 \end{cases}
 \label{eq:target-branches}
\end{equation}
In particular, $0\le\widetilde h\le10\rho^2$. The derivative bounds
needed below are
\begin{equation}
 0\le\rho\le2,\qquad
 \norm{\rho'}_\infty\le4,\qquad
 \Lip(\rho')\le40,
 \label{eq:rho-bounds}
\end{equation}
and $\rho=\rho'=0$ on $[-1/2,1/2]$. Define
\begin{equation}
 a_+(t)=5-\phi(t),\qquad a_-(t)=5+\phi(-t).
 \label{eq:a-branches}
\end{equation}
Both branches obey
\begin{equation}
 0<a_\pm(t)<10,\qquad |a_\pm'(t)|\le3,\qquad
 |a_\pm''(t)|\le3.
 \label{eq:a-bounds}
\end{equation}
Appendix~\ref{app:rho} verifies \eqref{eq:rho-bounds}.

The local block constructed in
Sections~\ref{sec:robust-chain}--\ref{sec:perspective} will have the
exact value
\begin{equation}
 \max_y C_V(s,t;y)=h(s,t).
 \label{eq:block-value}
\end{equation}
The next lemma records the consequences for assembly; its local
hypotheses are verified in Proposition~\ref{prop:local-block}.

\begin{lemma}[Assembly and primal compression]
\label{lem:outer-compat}
Suppose the local block satisfies \eqref{eq:block-value}, has a jointly
$L_C$-Lipschitz gradient with $L_C=4\cdot10^7$, and has primal-gradient
norm at most $2\cdot10^6$. Define
\begin{equation}
 \widehat f_{M,N,U,V}(x,y)
 =\sum_{i=1}^M C_{V_i}(s_{i-1},s_i;y^{(i)})
   +\frac1{10}\norm x^2,
 \quad s=U^\top\varrho(x),\quad s_0=1.
 \label{eq:rotated-saddle}
\end{equation}
Its value function is exactly \eqref{eq:outer-soft}, the primal
compression preserves every dual P\L{} inequality, and the full
function is jointly $\ell_0$-smooth with
\begin{equation}
 \ell_0=10^8.
 \label{eq:ell0}
\end{equation}
The initial value gap is at most $12M$, and
\eqref{eq:soft-obstruction} applies to this value function.
\end{lemma}

\begin{proof}
The dual blocks are disjoint, so their exact maxima add. Compression
acts only on $x$, and hence leaves the dual P\L{} inequality unchanged.
Also $\widehat\Phi(0)=F_M(0)$ and
$\inf\widehat\Phi\ge\inf F_M$, proving the gap bound.

For smoothness, color the path blocks by parity. Blocks of one color
have disjoint latent supports, so their sum has Hessian norm at most
$L_C$. The latent saddle therefore has Hessian norm at most
$2L_C=8\cdot10^7$. Each primal coordinate belongs to at most two
blocks, giving
$\norm{\nabla_s\sum_iC_{V_i}}\le4\cdot10^6\sqrt M$.
The compression has Jacobian norm at most one and a Jacobian
Lipschitz constant at most $3/R_x$. Its additional chain-rule
curvature is at most
\[
 \frac3{230\sqrt M}\,4\cdot10^6\sqrt M<5.3\cdot10^4.
\]
The public quadratic adds $1/5$. These bounds are below $10^8$.
The local blocks are $C^{1,1}$, and the same chain-rule estimates
bound gradient differences globally, so twice differentiability at the
scalar interfaces is not required. The gradient obstruction follows from
Lemma~\ref{lem:soft-projection}.
\end{proof}

\subsection{A dual chain with flat activation regions}
\label{sec:robust-chain}

The weighted chain follows Pan and Li~\cite{panli}. We shift both its
activation gate and its negative-part penalty so that an unrevealed
coordinate has a genuine interval of first-order silence. Set $\tau=1/4$
and define
\[
 q_0(t)=
 \begin{cases}
 0,&t\le0,\\
 3t^2-2t^3,&0<t<1,\\
 1,&t\ge1,
 \end{cases}
 \quad
 p(t)=q_0\!\left(\frac{t+1}{2}\right),
 \quad
 q_\tau(t)=q_0\!\left(\frac{t-\tau}{1-\tau}\right),
 \quad
 r_\tau(t)=(-t-\tau)_+.
 \label{eq:gates-def}
\]
The shift is applied both to the activation gate and to the negative-part penalty.  This is essential: shifting only $q$ leaves a first-order side channel in the nominal dead zone.

\begin{lemma}[Scalar gate estimates]
\label{lem:scalar-ledger}
For every $t\in\R$,
\begin{align}
 &0\le q_\tau(t)\le p(t)\le1,
 \qquad
 1-q_\tau(t)\le3\bigl(p'(t)^2+r_\tau(t)^2\bigr),
 \label{eq:scalar-main}\\
 &\norm{q_\tau'}_\infty\le2,
 \quad \Lip(q_\tau')\le\frac{32}{3},
 \quad \norm{p'}_\infty\le\frac34,
 \quad \Lip(p')\le\frac32.
 \label{eq:scalar-derivs}
\end{align}
Both $q_\tau$ and the derivative of $r_\tau^2/2$ vanish on $[-\tau,\tau]$.
\end{lemma}

\begin{proof}
The derivative bounds follow from $\norm{q_0'}_\infty=3/2$ and $\Lip(q_0')=6$.  For $t\le1$,
$(t-\tau)/(1-\tau)\le(t+1)/2$, and for $t\ge1$ both gates equal one; hence $q_\tau\le p$.
For \eqref{eq:scalar-main}, four intervals suffice.  On $[\tau,1)$, direct substitution gives
\[
 \frac{1-q_\tau(t)}{p'(t)^2}
 =\frac{256(8t+1)}{243(1+t)^2}\le\frac{4096}{1701}<3.
\]
For $t\ge1$, both sides vanish. On $[-\tau,\tau]$, $q_\tau=0$ and $p'(t)\ge45/64$, so $3p'(t)^2>1$.
On $(-1,-\tau)$, writing $s=-t$ gives
\[
 p'(t)^2+r_\tau(t)^2
 =\frac9{16}(1-s^2)^2+(s-1/4)^2\ge\frac{25}{72}>\frac13.
\]
Finally, for $t\le-1$, $r_\tau(t)^2\ge9/16$.  The flat-zone assertion is immediate.
\end{proof}

For $N\ge2$, define weights by
\begin{equation}
 \omega_{N-1}=\omega_N=1,
 \qquad
 \omega_{j-1}=\sqrt{\omega_j/N},
 \quad j=1,\ldots,N-1.
 \label{eq:weights}
\end{equation}
They obey
\begin{equation}
 \frac1N\le\omega_0\le\cdots\le\omega_N=1,
 \quad
 \sum_{j=1}^N\omega_{j-1}<3,
 \quad
 \sum_{j=1}^N\frac{\omega_{j-1}^2}{\omega_j}<2,
 \quad
 \sum_{j=1}^N\omega_j<4.
 \label{eq:weight-ledger}
\end{equation}
The choice of weights is due to Pan and Li~\cite{panli}; the estimates
can be checked directly. The recursion gives
$\omega_{N-1-k}=N^{-1+2^{-k}}$ for $0\le k\le N-1$, proving their
ordering. Since $N^t\le1+(N-1)t$ for $0\le t\le1$,
\[
 \sum_{j=1}^N\omega_{j-1}
 =\frac1N\sum_{k=0}^{N-1}N^{2^{-k}}
 \le1+\left(1-\frac1N\right)\sum_{k=0}^{N-1}2^{-k}<3.
\]
Moreover,
$\sum_{j=1}^N\omega_{j-1}^2/\omega_j=(N-1)/N+1<2$,
and adding $\omega_N$ while removing $\omega_0$ gives the last estimate.

Set $u_0=1$ and
\begin{align}
 H_\tau(u)
 &=\sum_{j=1}^N\left[
   \omega_{j-1}\bigl(1-q_\tau(u_{j-1})p(u_j)\bigr)
   +\frac{\omega_j}{2}r_\tau(u_j)^2
 \right],
 \label{eq:Htau}\\
 D_a^{\rm n}(u)
 &=H_\tau(u)+a\bigl(1-q_\tau(u_N)\bigr),
 \qquad 0\le a\le10.
 \label{eq:Dnatural}
\end{align}
Let $W=\diag(\sqrt{\omega_1},\ldots,\sqrt{\omega_N})$ and define the physical-coordinate deficit
\begin{equation}
 D_a(z):=D_a^{\rm n}(W^{-1}z).
 \label{eq:Dphysical}
\end{equation}

\begin{lemma}[Dual-chain estimates]
\label{lem:robust-chain}
For every $N\ge2$ and $a\in[0,10]$:
\begin{enumerate}
\item $D_a\ge0$, $D_a(W\one)=0$, and
\begin{equation}
 \norm{\nabla D_a(z)}^2\ge\frac{1}{385N}D_a(z).
 \label{eq:D-PL}
\end{equation}
\item $\nabla D_a$ is $128$-Lipschitz.
\item Every coordinate of $\nabla D_a$ is nonpositive, and $\partial_jD_a(z)=0$ for $z_j\ge\sqrt{\omega_j}$.
\item On $\norm z\le32$,
\begin{equation}
 0\le D_a(z)\le1040,
 \qquad
 \norm{\nabla D_a(z)}\le4097.
 \label{eq:D-ball}
\end{equation}
\item If $|u_j|\le\tau$ for all $j\ge k$, then no coordinate later than $k$ occurs in the gradient; the only nonconstant scalar term involving this suffix is its incoming $p(u_k)$ factor, whose derivative contains $p'(u_k)$.
\end{enumerate}
\end{lemma}

\begin{proof}
Write $\norm v_{*,\omega}^2=\sum_jv_j^2/\omega_j$.  Direct differentiation gives, for $j<N$,
\[
 \partial_jH_\tau
 =-\omega_{j-1}q_\tau(u_{j-1})p'(u_j)
  -\omega_jr_\tau(u_j)
  -\omega_jq_\tau'(u_j)p(u_{j+1}),
\]
and the same expression without the final term for $j=N$.  Every summand is nonpositive.  Hence
\begin{equation}
 \norm{\nabla H_\tau(u)}_{*,\omega}^2
 \ge\frac1N\sum_{j=1}^N
 \left[(q_\tau(u_{j-1})p'(u_j))^2+r_\tau(u_j)^2\right].
 \label{eq:H-grad-lower}
\end{equation}
Fix $k$.  If $q_\tau(u_k)<1$, choose the smallest $j\le k$ with
$1-q_\tau(u_j)\ge\tfrac12(1-q_\tau(u_k))$.  Then $q_\tau(u_{j-1})\ge1/2$, and Lemma~\ref{lem:scalar-ledger} gives
\begin{equation}
 1-q_\tau(u_k)
 \le24N\norm{\nabla H_\tau(u)}_{*,\omega}^2.
 \label{eq:activation-24}
\end{equation}
Also,
\[
 1-q_\tau(u_{j-1})p(u_j)
 \le1-q_\tau(u_{j-1})+1-q_\tau(u_j).
\]
Using \eqref{eq:weight-ledger} and \eqref{eq:activation-24}, the total gate part is at most
$144N\norm{\nabla H_\tau}_{*,\omega}^2$.  The penalty part is at most
$\tfrac N2\norm{\nabla H_\tau}_{*,\omega}^2$, and the terminal term is at most
$240N\norm{\nabla H_\tau}_{*,\omega}^2$.  Therefore
\[
 D_a^{\rm n}(u)
 \le385N\norm{\nabla H_\tau(u)}_{*,\omega}^2.
\]
The extra terminal derivative in $D_a^{\rm n}$ has the same sign as the terminal derivative of $H_\tau$, so its weighted gradient norm can only increase.  The change $z=Wu$ converts the weighted dual norm into the Euclidean norm and proves \eqref{eq:D-PL}.

For smoothness, each gate has diagonal physical second derivatives bounded by
$32/3$ and $3/2$, and mixed derivative bounded by $3/2$.  The shifted quadratic penalty contributes at most one.  The only larger term is the terminal $a$-term, whose physical second derivative is at most $10(32/3)$.  Gershgorin's bound gives an almost-everywhere Hessian norm below $128$; continuity of the first derivative gives global $128$-smoothness.

The sign and saturation claims follow from monotonicity and the flat portions of the gates.  For the ball bound, use
\[
 H_\tau(u)\le3+\norm z^2+\tau^2\sum_j\omega_j
 \le\norm z^2+\frac{13}{4},
\]
then add the terminal term $a\le10$.  At the origin only the first physical coordinate of the gradient can be nonzero, with magnitude at most $3/4$; smoothness then gives
$\norm{\nabla D_a(z)}\le3/4+128\norm z\le4097$ on the radius-$32$ ball.  The final assertion follows because both $q_\tau$ and the derivative of the shifted penalty vanish on $[-\tau,\tau]$.
\end{proof}

\subsection{A radial cap that preserves the dual P\L{} inequality}
\label{sec:cap}

Set
\begin{equation}
 R=16,
 \qquad A=1,
 \qquad C=32,
 \qquad T_a(r)=a+C+Ar^2.
 \label{eq:cap-constants}
\end{equation}
Let
\[
 \chi(r)=
 \begin{cases}
 1,&r\le R,\\
 1-q_0((r-R)/R),&R<r<2R,\\
 0,&r\ge2R.
 \end{cases}
\]
On the annulus,
\begin{equation}
 R^2\chi'(r)^2\ge4\min\{\chi(r),1-\chi(r)\},
 \quad
 |\chi'|\le\frac{3}{2R},
 \quad
 \norm{\nabla^2\chi(\norm\cdot)}\le\frac6{R^2}.
 \label{eq:chi-ledger}
\end{equation}
Let $V\in\R^{d\times N}$ have orthonormal columns.  Define
\begin{equation}
 \mathscr D_{a,V}(w)
 =\chi(\norm w)D_a(V^\top w)
 +(1-\chi(\norm w))T_a(\norm w).
 \label{eq:capped-deficit}
\end{equation}

\begin{lemma}[Capped deficit]
\label{lem:cap}
Uniformly over $N\ge2$, $d\ge N$, $a\in[0,10]$, and $V$,
\begin{align}
 &\norm{\nabla\mathscr D_{a,V}(w)}^2
 \ge\frac{10^{-5}}{N}\mathscr D_{a,V}(w),
 \label{eq:cap-pl}\\
 &\Lip(\nabla\mathscr D_{a,V})\le1024.
 \label{eq:cap-smooth}
\end{align}
The minimum is zero and is attained at $w=VW\one$.  For $\norm w\ge32$,
\begin{equation}
 \mathscr D_{a,V}(w)=a+32+\norm w^2,
 \label{eq:public-tail}
\end{equation}
which is independent of the hidden frame.
\end{lemma}

\begin{proof}
On $\norm w\le16$ this is Lemma~\ref{lem:robust-chain}; on $\norm w\ge32$ it is a radial quadratic.  Consider $16<r:=\norm w<32$, set $z=V^\top w$, $e=w/r$, and $g=V\nabla D_a(z)$.  The sharper estimate obtained in the proof of Lemma~\ref{lem:robust-chain},
$D_a(z)\le\norm z^2+53/4$, gives
\[
 T_a(r)-D_a(z)
 \ge a+32+r^2-\norm z^2-\frac{53}{4}
 \ge32-\frac{53}{4}>18.
\]
Differentiation gives
\begin{equation}
 \nabla\mathscr D_{a,V}=\chi g+be,
 \qquad
 b=2r(1-\chi)+(-\chi')(T_a-D_a)\ge0.
 \label{eq:cap-gradient-decomp}
\end{equation}
A negative component of $\nabla D_a$ can oppose the outward radial direction only when
$0<z_j<\sqrt{\omega_j}$.  Since $\sum_j\omega_j<4$ and $r\ge16$,
\[
 \ip{g}{e}\ge-\frac18\norm g.
\]
Therefore
\begin{equation}
 \norm{\chi g+be}^2
 \ge\frac78\bigl(\chi^2\norm g^2+b^2\bigr).
 \label{eq:cap-no-cancel}
\end{equation}
If $\chi\ge1/2$, then
$\chi^2\norm g^2\ge\chi D_a/(770N)$.  Also
$T_a\le1066$ on the annulus and \eqref{eq:chi-ledger} gives
$b^2\ge5(1-\chi)\ge(1-\chi)T_a/214$.  If $\chi<1/2$, then
$b\ge2r(1-\chi)\ge16$, while
$\mathscr D_{a,V}\le1066$.  In both cases \eqref{eq:cap-pl} follows with the deliberately conservative constant $10^{-5}$.

For smoothness on the annulus, use
\[
 \norm{\nabla^2D_a}\le128,
 \quad \norm{\nabla D_a}\le4097,
 \quad D_a\le1040,
 \quad \norm{\nabla T_a}\le64,
 \quad T_a\le1066.
\]
Twice differentiating \eqref{eq:capped-deficit} and using \eqref{eq:chi-ledger} gives the operator-norm bound
\[
 128+2+2\frac{3}{32}(4097+64)+\frac6{16^2}(1040+1066)<1024.
\]
The value and first derivative match at both annulus boundaries, hence the global gradient is $1024$-Lipschitz.  Finally $\norm{W\one}<2<16$, so the zero of $D_a$ lies in the uncapped region.
\end{proof}

\subsection{The capped perspective block}
\label{sec:perspective}

Set
\begin{equation}
 Q_{a,V}(w)=a-\mathscr D_{a,V}(w).
 \label{eq:Qdef}
\end{equation}
On $\norm w\le32$, Lemma~\ref{lem:cap} and the preceding bounds imply
\begin{equation}
 |Q_{a,V}(w)|\le1100,
 \quad \norm{\nabla_wQ_{a,V}(w)}\le4500,
 \quad \norm{\nabla^2_{ww}Q_{a,V}(w)}\le1024.
 \label{eq:Q-spatial-bounds}
\end{equation}
The parameter dependence is affine and
\begin{equation}
 \partial_aQ_{a,V}(w)
 =\chi(\norm w)q_\tau(v_N^\top w),
 \label{eq:Qa-formula}
\end{equation}
so on the same ball
\begin{equation}
 |Q_a|\le1,
 \qquad \norm{\nabla_wQ_a}\le3,
 \qquad \norm{\nabla^2_{ww}Q_a}\le12.
 \label{eq:Qa-bounds}
\end{equation}
For $\norm w\ge32$, one has exactly
\begin{equation}
 Q_{a,V}(w)=-32-\norm w^2,
 \label{eq:Q-tail}
\end{equation}
independently of $a$ and $V$.

For $\eta>0$ and $0\le\alpha\le10\eta^2$, define
\begin{equation}
 \mathscr B_V(\eta,\alpha;y)
 =\eta^2Q_{\alpha/\eta^2,V}(y/\eta),
 \label{eq:Bcap}
\end{equation}
and set $\mathscr B_V(0,0;y)=-\norm y^2$.

\begin{lemma}[Global perspective extension]
\label{lem:perspective-extension}
Let $a(t)$ be either branch $a_+(t)$ or $a_-(t)$ from \eqref{eq:a-branches}.  The function
\begin{equation}
 P_{V,a}(\eta,t,y)=
 \begin{cases}
 \eta^2Q_{a(t),V}(y/\eta),&\eta>0,\\
 -\norm y^2,&\eta=0,
 \end{cases}
 \label{eq:Pdef}
\end{equation}
has globally Lipschitz gradient on $[0,2]\times\R\times\R^d$.  In the hidden cone
$\norm y\le32\eta$, its almost-everywhere second derivatives obey
\begin{align}
 |P_{\eta\eta}|&\le1.4\cdot10^6,
 &\norm{P_{y\eta}}&\le3.8\cdot10^4,
 &\norm{P_{yy}}&\le1024,
 \label{eq:P-hess-1}\\
 |P_{\eta t}|&\le588,
 &\norm{P_{yt}}&\le18,
 &|P_{tt}|&\le12.
 \label{eq:P-hess-2}
\end{align}
Outside the hidden cone, $P_{V,a}=-32\eta^2-\norm y^2$ exactly.
\end{lemma}

\begin{proof}
For $w=y/\eta$ and $\eta>0$,
\begin{align*}
 P_\eta&=\eta(2Q-\ip{w}{Q_w}),
 &P_y&=\eta Q_w,\\
 P_{\eta\eta}&=2Q-2\ip{w}{Q_w}+w^\top Q_{ww}w,
 &P_{y\eta}&=Q_w-Q_{ww}w,
 &P_{yy}&=Q_{ww},\\
 P_{\eta t}&=\eta(2Q_a-\ip{w}{Q_{aw}})a'(t),
 &P_{yt}&=\eta Q_{aw}a'(t),
 &P_{tt}&=\eta^2Q_a a''(t).
\end{align*}
The bounds follow from \eqref{eq:Q-spatial-bounds}--\eqref{eq:Qa-bounds},
$\norm w\le32$, $\eta\le2$, and \eqref{eq:a-bounds}.  In the complement of the cone,
\eqref{eq:Q-tail} gives the stated public quadratic.  For every positive lower bound on $\eta$, the formula is locally
$C^{1,1}$, so its gradient is absolutely continuous along line segments.
The displayed Hessian estimates give a Lipschitz bound independent of
that lower bound. In the cone, $\norm y\le32\eta$, all first derivatives
converge to their prescribed values as $\eta\downarrow0$; outside the
cone this follows from the exact quadratic formula. Approximating
boundary segments by segments with $\eta>0$ proves the same Lipschitz
estimate up to $\eta=0$.
\end{proof}

Define the capped local saddle block directly from the target parameter:
\begin{equation}
 C_V(s,t;y)
 =-5\rho(s)^2
  +\mathscr B_V\bigl(\rho(s),\widetilde h(s,t);y\bigr),
 \label{eq:local-block}
\end{equation}
where the value at \(\rho(s)=0\) is the extension
\(\mathscr B_V(0,0;y)=-\norm y^2\).  On the two nonzero branches this is
\(-5\rho(s)^2+P_{V,a_+}(\rho(s),t,y)\) and
\(-5\rho(s)^2+P_{V,a_-}(\rho(s),t,y)\), respectively.  Thus the sign branch is part of the definition and is not inferred from the symmetric scale \(\rho(s)\).

\begin{proposition}[Local saddle-block properties]
\label{prop:local-block}
For all $s,t,y$ and every hidden frame $V$:
\begin{enumerate}
\item
\begin{equation}
 \max_yC_V(s,t;y)=h(s,t),
 \label{eq:local-value}
\end{equation}
and a maximizer is $y=\rho(s)VW\one$.
\item
\begin{equation}
 \norm{\nabla_yC_V(s,t;y)}^2
 \ge\frac{10^{-5}}{N}\bigl(h(s,t)-C_V(s,t;y)\bigr).
 \label{eq:local-dual-pl}
\end{equation}
Thus, in the convention \eqref{eq:model-pl}, the unscaled dual P\L{} modulus is
\begin{equation}
 \frac{\mu_0}{N},
 \qquad \mu_0=5\cdot10^{-6}.
 \label{eq:mu0}
\end{equation}
\item $C_V$ has jointly $L_C$-Lipschitz gradient with
\begin{equation}
 L_C=4\cdot10^7.
 \label{eq:LC}
\end{equation}
Its primal-gradient norm is at most $2\cdot10^6$.
\item If $\rho(s)=0$, then $C_V(s,t;y)=-\norm y^2$, independent of $t$ and $V$.
\item For $\rho(s)>0$,
\begin{equation}
 \partial_{\widetilde h}\mathscr B_V
 =\chi(\norm y/\rho(s))
 q_\tau\!\left(
  \frac{v_N^\top y}{\rho(s)\sqrt{\omega_N}}
 \right).
 \label{eq:successor-gate}
\end{equation}
Hence the successor primal variable is absent from the gradient whenever the terminal normalized dual coordinate belongs to $[-\tau,\tau]$, and also throughout the public tail.
\end{enumerate}
\end{proposition}

\begin{proof}
On the positive branch set $a_\star(t)=a_+(t)$ and on the negative branch set $a_\star(t)=a_-(t)$.  The gap is
\[
 h(s,t)-C_V(s,t;y)
 =\rho(s)^2\mathscr D_{a_\star(t),V}(y/\rho(s)),
\]
when $\rho(s)>0$. Lemma~\ref{lem:cap} proves the value and P\L{}
claims in this case. At $\rho(s)=0$, the gap is $\norm y^2$ and
$\norm{\nabla_yC_V}^2=4\norm y^2$, so the same inequality holds
directly, with maximizer $y=0$. Formula~\eqref{eq:Qa-formula} gives
\eqref{eq:successor-gate}.

For smoothness, apply Lemma~\ref{lem:perspective-extension} on each sign branch and substitute \eqref{eq:rho-bounds} into the chain rule for
$-5\rho^2+P_{V,a_\star}(\rho,t,y)$.  The largest entry is
\[
 |C_{ss}|
 \le(1.4\cdot10^6+10)\,4^2
 +(3.0\cdot10^5+20)\,40
 <3.5\cdot10^7.
\]
The remaining mixed entries are below $1.6\cdot10^5$.  Inside $\norm y\le32\rho(s)$, we used the first-derivative consequences of \eqref{eq:Q-spatial-bounds},
\[
 |P_\eta|\le3.0\cdot10^5,
 \qquad
 \norm{P_y}\le9.0\cdot10^3,
 \qquad
 |P_t|\le12.
\]
Outside this cone, $P=-32\eta^2-\norm y^2$, so
$P_\eta=-64\eta$, $P_t=0$, and $P_{yy}=-2I$. A block row-sum bound
therefore gives $L_C=4\cdot10^7$ globally, and the same estimates bound
the primal gradient. At $s=\pm1/2$, the identities $\rho=\rho'=0$
make the value and first derivatives agree with the public quadratic.
The uniform branchwise Lipschitz bounds consequently extend across
both interfaces.
\end{proof}

Combining Proposition~\ref{prop:local-block} with Lemma~\ref{lem:outer-compat} gives uniform constants for the complete unscaled construction:
\begin{equation}
 \boxed{
 \ell_0=10^8,
 \quad \mu_0=5\cdot10^{-6},
 \quad c_{\rm cap}=10^{-5},
 \quad L_C=4\cdot10^7.
 }
 \label{eq:common-ledger}
\end{equation}
These conservative choices are independent of both chain lengths and all ambient dimensions.

\subsection{From the value gradient to the weighted joint residual}
\label{sec:transfer}

For block $i$, define the nonnegative local deficit
\[
 d_i(s_{i-1},s_i;y^{(i)})
 =h(s_{i-1},s_i)-C_{V_i}(s_{i-1},s_i;y^{(i)}).
\]
For fixed $y^{(i)}$, the deficit is nonnegative as a function of its
two primal arguments. The bounds in \eqref{eq:rho-bounds} and
\eqref{eq:a-bounds} give $\Lip(\nabla h)\le2000$; together with
Proposition~\ref{prop:local-block}, they allow the primal smoothness
constant
\begin{equation}
 L_{\rm def}=5\cdot10^7.
 \label{eq:Ldef}
\end{equation}
The self-bounding inequality for a nonnegative $L_{\rm def}$-smooth function gives
$\norm{\nabla_{s,t}d_i}^2\le2L_{\rm def}d_i$.

\begin{lemma}[Local-deficit comparison]
\label{lem:joint-transfer}
For the soft-projected saddle construction \eqref{eq:rotated-saddle},
\begin{equation}
 \norm{\nabla\widehat\Phi(x)-\nabla_x\widehat f(x,y)}^2
 \le4L_{\rm def}\bigl(\widehat\Phi(x)-\widehat f(x,y)\bigr)
 \le K_{\rm tr}^2N\norm{\nabla_y\widehat f(x,y)}^2,
 \label{eq:transfer-chain}
\end{equation}
where one may take
\begin{equation}
 K_{\rm tr}^2=\frac{4L_{\rm def}}{10^{-5}}=2\cdot10^{13}.
 \label{eq:Ktr}
\end{equation}
If
\begin{equation}
 N\le\frac{\kappa}{16K_{\rm tr}^2},
 \label{eq:N-transfer-condition}
\end{equation}
then
\begin{equation}
 \norm{\nabla_x\widehat f(x,y)}^2
 +\kappa\norm{\nabla_y\widehat f(x,y)}^2
 \ge\frac14\norm{\nabla\widehat\Phi(x)}^2.
 \label{eq:joint-transfer-final}
\end{equation}
\end{lemma}

\begin{proof}
Each latent primal coordinate occurs in at most two neighboring deficits.  Hence
\[
 \norm{\nabla_sF_M(s)-\nabla_s\bar f(s,y)}^2
 \le2\sum_i\norm{\nabla_{s,t}d_i}^2
 \le4L_{\rm def}\sum_i d_i.
\]
The Jacobian of the primal compression and the frame $U$ are contractions; the public quadratic cancels from the gradient difference.  This proves the first inequality.  Summing \eqref{eq:local-dual-pl} over the disjoint dual blocks gives
$\sum_id_i\le10^5N\norm{\nabla_y\widehat f}^2$, proving the second.
Let $a=\norm{\nabla_x\widehat f}$ and
$b=\sqrt\kappa\norm{\nabla_y\widehat f}$.  Under \eqref{eq:N-transfer-condition},
$\norm{\nabla\widehat\Phi}\le a+b/4$.  Cauchy--Schwarz yields
$(a+b/4)^2\le(17/16)(a^2+b^2)$, which is stronger than \eqref{eq:joint-transfer-final}.
\end{proof}

In particular, the low-progress value-gradient obstruction \eqref{eq:soft-obstruction} implies the unscaled pointwise obstruction
\begin{equation}
 \norm{\nabla_x\widehat f(x,y)}^2
 +\kappa\norm{\nabla_y\widehat f(x,y)}^2
 \ge\frac1{16}.
 \label{eq:unscaled-joint-obstruction}
\end{equation}

\section{Information revealed by one oracle response}
\label{sec:prefix}

The construction imposes an order on the latent coordinates. This order
is used only in the proof; admissible algorithms are not required to
respect it. The interleaved order is
\begin{equation}
 y^{(1)}_1,\ldots,y^{(1)}_N,x_1,
 y^{(2)}_1,\ldots,y^{(2)}_N,x_2,\ldots,
 y^{(M)}_1,\ldots,y^{(M)}_N,x_M.
 \label{eq:global-order}
\end{equation}
Write $K=M(N+1)$ for its length and let $\pi_r$ denote the $r$th position.  A dual position $(i,j)$ is associated with the hidden column $v_{i,j}$ of $V_i$; a primal connector $i$ is associated with the hidden column $u_i$ of $U$.

For an ambient query $z=(x,y)$, put $s=U^\top\varrho(x)$.  In block $i$ with $\eta_i=\rho(s_{i-1})>0$, define
\[
 w_i=y^{(i)}/\eta_i,
 \qquad
 u_{i,j}=\frac{v_{i,j}^\top y^{(i)}}{\eta_i\sqrt{\omega_j}}.
\]
A block with $\eta_i=0$ is public.  A block with $\norm{w_i}\ge32$ is also public by \eqref{eq:public-tail}.  Public blocks are ignored by the hidden progress statistic.

The robust dead zones are
\begin{equation}
 \begin{aligned}
  |s_i|&\le\frac14
  &&\text{for an unrevealed primal position},\\
  |u_{i,j}|&\le\tau=\frac14
  &&\text{for an unrevealed dual position}.
 \end{aligned}
 \label{eq:robust-dead-zones}
\end{equation}
For the concentration argument we use the stronger, weight-adapted dual threshold
\begin{equation}
 |v_{i,j}^\top y^{(i)}|
 \le c_*\eta_i\sqrt{\omega_{j-1}},
 \qquad c_*:=\frac18.
 \label{eq:weight-threshold}
\end{equation}
Since $\omega_{j-1}/\omega_j\le1$, this implies $|u_{i,j}|\le1/8<\tau$.

\subsection{Prefix-truncated functions}

A conditional-Haar argument requires every ideal response to be a
function of revealed columns before conditioning on a favorable
projection event. We enforce this requirement by truncating latent
coordinates in the scalar objective before differentiating.

For $r\in\{0,1,\ldots,K\}$, define a prefix-truncated latent representation as follows.
For every primal coordinate $i$, retain $s_i=u_i^\top\varrho(x)$ if the connector position $(i,N+1)$ is at most $r$, and replace it by zero otherwise.  For every dual coordinate $(i,j)$, retain $v_{i,j}^\top y^{(i)}$ if its position is at most $r$, and replace it by zero otherwise.  The public radial quantity $\norm{y^{(i)}}$ is never truncated.  Substituting these values into the capped saddle construction defines a globally $C^{1,1}$ scalar function
\begin{equation}
 \widehat f^{[r]}(x,y).
 \label{eq:prefix-function}
\end{equation}
It depends only on the public constants and on hidden frame columns whose positions are at most $r$.  Set
\begin{equation}
 \Delta_r(z)=\nabla\widehat f^{[r]}(z)-\nabla\widehat f^{[r-1]}(z),
 \qquad r=1,\ldots,K.
 \label{eq:prefix-increment}
\end{equation}

\begin{lemma}[Prefix identity]
\label{lem:prefix-identity}
Fix $0\le r<K$ and a query $z$.  Suppose that every primal position
after $r$ satisfies $|s_i|\le1/4$, and that every nonpublic dual position
after $r$ satisfies
\begin{equation}
 |u_{i,j}|\le\tau.
 \label{eq:prefix-dead-zone}
\end{equation}
Then
\begin{equation}
 \nabla\widehat f(z)=\nabla\widehat f^{[r+1]}(z).
 \label{eq:prefix-equality}
\end{equation}
If position $r+1$ belongs to a public dual block, then this gradient also
equals $\nabla\widehat f^{[r]}(z)$.  For every query, without any
projection hypothesis, $\nabla\widehat f^{[r]}(z)$ depends only on $z$
and the frame columns in positions at most $r$.
\end{lemma}

\begin{proof}
The last assertion follows directly from the definition of the prefix
function.  To prove the identity, first consider a dual frontier
$(i,k)$.  Within this block, all normalized coordinates from $k$ onward
lie in $[-\tau,\tau]$.  Their activation gates and gate derivatives,
and their shifted-penalty derivatives, vanish.  Consequently, the only
possibly nonconstant term involving this suffix is the incoming factor
$p(u_{i,k})$, which is retained in prefix $r+1$.  The terminal activation
factor vanishes, so neither the value nor its first derivatives depend
on the unrevealed connector $s_i$.  Since every later predecessor
primal coordinate lies in $[-1/4,1/4]$, every later block has
$\rho=\rho'=0$ and is the public quadratic.  If block $i$ is itself
public, it has no hidden dependence either, and prefix $r$ suffices.

If the frontier is a primal connector $i$, its entire predecessor dual
block is retained, and prefix $r+1$ also retains this connector.  All
later blocks have predecessor scale and scale derivative zero.  Thus
truncating every later latent projection changes no value or
first-derivative term.  The radial norms and the public primal
quadratic are never truncated, so the same conclusion holds after all
chain-rule terms are included.
\end{proof}

\subsection{The three boundary cases}

\begin{lemma}[Separation of adjacent blocks]
\label{lem:boundary-audit}
On the event in Lemma~\ref{lem:prefix-identity}, one oracle response can introduce at most the current position in \eqref{eq:global-order}.  More precisely:
\begin{enumerate}
\item If the current position is $y^{(i)}_j$ with $j<N$, the only new frame column on which the full increment can depend is
$v_{i,j}$. It may also have components along previously revealed columns
and public radial vectors. The next column $v_{i,j+1}$ is absent because
$q_\tau(u_{i,j})=q_\tau'(u_{i,j})=0$ at the current coordinate.
\item If the current position is $y^{(i)}_N$, an exact response may depend on $v_{i,N}$, but it can be computed without the unrevealed connector column
$u_i$ or any later column: its latent derivative with respect to $s_i$
vanishes by \eqref{eq:successor-gate}.
\item If the current position is the connector $x_i$, an exact response may depend on $u_i$, but it contains no column of $V_{i+1}$: before $u_i$ is revealed, $|s_i|\le1/4$, so $\rho(s_i)=\rho'(s_i)=0$ and block $i+1$ is public.
\end{enumerate}
\end{lemma}

\begin{proof}
The first claim is the flat-zone statement in Lemma~\ref{lem:robust-chain}.  In the terminal case, $|u_{i,N}|\le\tau$, so both $q_\tau(u_{i,N})$ and $q_\tau'(u_{i,N})$ vanish; the target derivative with respect to $s_i$ is therefore zero.  In the connector case, the outer flat zone and all of its derivatives are zero on $[-1/2,1/2]$, and the stronger threshold $1/4$ lies strictly inside it.
\end{proof}

These statements concern dependence on hidden frame columns, not
orthogonality of the ambient response: public radial terms may have
nonzero projections onto unrevealed directions. The identities prevent
a terminal dual reveal and its successor connector, or a connector and
the next dual coordinate, from being revealed in the same response.

\subsection{Bounding the complete frontier increment}
\label{sec:increment}

Let position $r$ be an inner dual frontier $(i,k)$ with $k<N$.
At a query where its normalized coordinate lies in the flat region,
the only nonconstant scalar difference between prefix $r$ and prefix
$r-1$ is
\begin{equation}
 R_k(\eta,y)
 =\eta^2\chi(\norm y/\eta)\,
 \omega_{k-1}q_\tau(u_{k-1})
 \bigl(p(u_k)-p(0)\bigr).
 \label{eq:Rk}
\end{equation}
It is important to mask the gradient of the entire scalar $R_k$, not only its component along $v_{i,k}$.

\begin{lemma}[Uniform whole-gradient frontier bound]
\label{lem:frontier-bound}
Assume $\eta>0$, $\norm y\le32\eta$, and that the current coordinate satisfies
$|v_{i,k}^\top y|\le c_*\eta\sqrt{\omega_{k-1}}$.  Then the difference between the complete ambient gradients of the two neighboring prefix functions obeys
\begin{equation}
 \norm{\Delta_r(z)}\le\frac{8}{\sqrt N}.
 \label{eq:increment-unscaled}
\end{equation}
The vector $\nabla\widehat f^{[r-1]}(z)$ is a measurable function of
the query and the previously revealed columns only. The increment
$\Delta_r(z)$ additionally uses the current column, but no later column.
\end{lemma}

\begin{proof}
Put
\[
 \alpha_k:=\sqrt{\frac{\omega_{k-1}}{\omega_k}}\le1,
 \qquad
 \beta_k:=\frac{\omega_{k-1}}{\sqrt{\omega_k}}=\frac1{\sqrt N}.
\]
The threshold gives $|u_k|\le\alpha_k/8$, and hence
\begin{equation}
 |p(u_k)-p(0)|
 \le\frac34|u_k|
 \le\frac{3\alpha_k}{32}.
 \label{eq:delta-p-bound}
\end{equation}
Differentiate the scalar $R_k$ in \eqref{eq:Rk}.  Its $y$-gradient is the sum of the direct current-column term, the predecessor-gate term, and the radial-cutoff term.  Their norms are respectively at most
\begin{align*}
 &\eta\frac{\omega_{k-1}}{\sqrt{\omega_k}}\norm{p'}_\infty
   \le\frac{3\eta}{4\sqrt N},\\
 &\eta\sqrt{\omega_{k-1}}\norm{q_\tau'}_\infty
   |p(u_k)-p(0)|
   \le\frac{3\eta}{16\sqrt N},\\
 &\eta\norm{\chi'}_\infty\omega_{k-1}|p(u_k)-p(0)|
   \le\frac{9\eta}{1024\sqrt N}.
\end{align*}
For the derivative with respect to the perspective scale $\eta$, the prefactor, radial cutoff, predecessor normalization, and current normalization contribute at most
\[
 \frac{3\eta}{16\sqrt N},
 \qquad
 \frac{9\eta}{32\sqrt N},
 \qquad
 \frac{3\eta}{16\sqrt N},
 \qquad
 \frac{3\eta}{32\sqrt N},
\]
respectively.  In the predecessor term we used that
$q_\tau'(u_{k-1})\ne0$ implies $|u_{k-1}|\le1$; the term is absent for $k=1$ because $u_0\equiv1$.
Thus
\[
 \norm{\nabla_yR_k}<\frac{\eta}{\sqrt N},
 \qquad
 |\partial_\eta R_k|\le\frac{3\eta}{4\sqrt N}.
\]
Composition with $\eta=\rho(s_{i-1})$ multiplies the second bound by at most
$\norm{\rho'}_\infty\le4$, and the primal frame and soft projection are contractions.  Since $\eta\le2$, the sum of the ambient primal and dual gradient norms is below $8/\sqrt N$, proving \eqref{eq:increment-unscaled}.

All later latent coordinates are set to zero in both prefix functions.
The prefix definition removes the current column from the base before
differentiation. Thus the base uses only earlier columns, and the
increment uses the current column but no later one.
\end{proof}

\section{Scaling and initialization}
\label{sec:scaling}

Define
\begin{equation}
 c_N:=\min\left\{
  \frac{\mu_0}{4\ell_0},
  \frac{1}{16K_{\rm tr}^2}
 \right\},
 \qquad
 c_M:=\frac{1}{12800\ell_0}.
 \label{eq:cN-cM}
\end{equation}
With the constants \eqref{eq:common-ledger} and \eqref{eq:Ktr}, these are positive absolute constants; numerically $c_N\le3.125\cdot10^{-15}$.  Set
\begin{equation}
 N=\lfloor c_N\kappa\rfloor,
 \qquad
 M=\left\lfloor c_M\frac{LH}{\eps^2}\right\rfloor.
 \label{eq:MN}
\end{equation}
One explicit choice of regime constants is
\begin{equation}
 \kappa_0=\frac8{c_N},\qquad c_\eps=\frac{c_M}{8}.
 \label{eq:regime-constants}
\end{equation}
The hypotheses of Theorem~\ref{thm:main-lower} then give $N,M\ge8$.
In particular, $S=M(N-1)\ge56$, so rounding the query horizon only
changes a universal constant.

Let
\begin{equation}
 \lambda=8\ell_0\frac{\eps}{L},
 \qquad
 f(x,y)=\frac{L\lambda^2}{\ell_0}
 \widehat f_{M,N,U,V}(x/\lambda,y/\lambda).
 \label{eq:scaled-function}
\end{equation}
The same formula defines every scaled prefix function $f^{[r]}$.

\begin{lemma}[Admissibility after scaling]
\label{lem:scaled-admissible}
The function \eqref{eq:scaled-function} is $L$-smooth, satisfies the dual P\L{} condition with modulus at least $L/\kappa$, has finite attained inner maxima, and obeys the charged initialization \eqref{eq:charged-init} at $z_0=x_0=y_0=0$.
\end{lemma}

\begin{proof}
Smoothness follows from \eqref{eq:ell0}.  The scaled dual modulus is
\[
 \frac{L\mu_0}{\ell_0N}
 \ge\frac{L}{\kappa}
\]
by the first choice in \eqref{eq:cN-cM}.  Attainment follows from the explicit maximizer of each capped block.

The unscaled value gap is at most $12M$.  At $y=0$, the target parameter cancels and
$C_V(s,t;0)=-(5+H_\tau(0))\rho(s)^2\ge-32$,
since $H_\tau(0)<3$ and $\rho^2\le4$. Hence
\[
 \inf_x\left[
   \widehat f_{M,N,U,V}(x,0)+\ell_0\norm x^2
 \right]\ge-32M.
\]
At $x=0$, all latent primal coordinates except the fixed $s_0=1$ vanish.  Since
$|h(s,t)|\le5\rho(s)^2\le20$ and $h(0,0)=0$, the initial value obeys
$|\widehat\Phi(0)|\le20$.  Denoting the auxiliary gap at the unscaled parameters by $\widehat b_0$, we obtain
\[
 \widehat b_0
 \le2(20+32M),
 \qquad
 \widehat\Phi(0)-\inf\widehat\Phi+\widehat b_0
 \le12M+40+64M\le96M
\]
for $M\ge2$.  We use the round bound $100M$.  Under the scaling, both the value gap and the auxiliary gap are multiplied by
\[
 \frac{L\lambda^2}{\ell_0}
 =64\ell_0\frac{\eps^2}{L}.
\]
By \eqref{eq:MN} and the definition of $c_M$, their sum is at most $H/2$ (with room for floors), proving \eqref{eq:charged-init}.
\end{proof}

\begin{lemma}[Scaled low-progress obstruction and stochastic signal]
\label{lem:scaled-obstruction}
At every point $(x,y)$ satisfying
$\prog_{1/4}(U^\top\varrho(x/\lambda))<M$,
\begin{equation}
 \norm{\nabla_x f(x,y)}^2
 +\kappa\norm{\nabla_y f(x,y)}^2\ge4\eps^2.
 \label{eq:scaled-joint-obstruction}
\end{equation}
If position $r$ is an inner dual frontier and its unscaled query
satisfies the hypotheses of Lemma~\ref{lem:frontier-bound}, then
$\norm{\nabla f^{[r]}-\nabla f^{[r-1]}}\le64\eps/\sqrt N$.
\end{lemma}

\begin{proof}
The gradient multiplier in \eqref{eq:scaled-function} is
$L\lambda/\ell_0=8\eps$.  At the stated low-progress point,
\eqref{eq:unscaled-joint-obstruction} applies to $(x/\lambda,y/\lambda)$.
Multiplication by $(8\eps)^2$ proves
\eqref{eq:scaled-joint-obstruction}.  Under the stated frontier hypotheses,
Lemma~\ref{lem:frontier-bound} gives a scaled increment norm at most
$8\eps\cdot8/\sqrt N=64\eps/\sqrt N$.
\end{proof}

\section{A fixed globally valid stochastic first-order oracle}
\label{sec:oracle}

The oracle must satisfy its moment conditions at every query, including queries
outside the region used in the lower-bound argument.  We therefore use the
prefix construction only after a deterministic test of the query and return
the exact gradient otherwise.  The test is part of a single fixed oracle map;
it does not depend on the transcript.

Let $f$ be the scaled construction in \eqref{eq:scaled-function}, and write
$\bar z=(\bar x,\bar y)=(x/\lambda,y/\lambda)$ for the corresponding unscaled
query.  Every progress statistic, cone condition, and projection threshold in
this section is evaluated at $\bar z$.  In particular,
\[
 s=U^\top\varrho(\bar x),\qquad
 \eta_i=\rho(s_{i-1}),\qquad
 u_{i,j}=\frac{v_{i,j}^\top\bar y^{(i)}}{\eta_i\sqrt{\omega_j}}
 \quad(\eta_i>0).
\]
Define $r=\rprog(\bar z)$ to be the largest position in
\eqref{eq:global-order} outside its robust dead zone, ignoring dual positions
in public blocks; set $r=0$ if there is no such position.  Thus every later primal position and every later nonpublic dual
position lies in its robust dead zone.  The simultaneous test asks whether
every later primal position satisfies $|s_i|\le1/4$ and every later nonpublic
dual position satisfies \eqref{eq:weight-threshold}, with $y^{(i)}$ there
replaced by $\bar y^{(i)}$.

Set
\begin{equation}
 \Gamma=\frac{64\eps}{\sqrt N},\qquad
 p=\frac{\Gamma^2}{\Gamma^2+\sigma^2},
 \label{eq:p-Gamma}
\end{equation}
and draw an independent $B\sim\operatorname{Bernoulli}(p)$ on each call.
If the simultaneous test passes and $r+1$ is an inner dual position
$(i,k)$ with $k<N$ in a nonpublic block, return
\begin{equation}
 G(z;B)=\nabla f^{[r]}(z)
   +\frac{B}{p}\bigl(\nabla f^{[r+1]}(z)-\nabla f^{[r]}(z)\bigr).
 \label{eq:fixed-oracle}
\end{equation}
In all other cases, return $\nabla f(z)$.  These exact cases include a failed
test, a public frontier, a terminal dual frontier, a primal connector, and
$r=K$, when there is no remaining frontier.  The Bernoulli variable is unused
on an exact branch.  All quantities defining this map are measurable
functions of $z$ and the fixed instance.

\begin{lemma}[Global oracle validity]
\label{lem:oracle-valid}
For every fixed hard instance and every query $z$,
\[
 \E_B[G(z;B)]=\nabla f(z),\qquad
 \E_B\norm{G(z;B)-\nabla f(z)}^2\le\sigma^2.
\]
The same identities hold conditionally on the past at every adaptive call.
\end{lemma}

\begin{proof}
There is nothing to prove on an exact branch.  On a masked branch, the
prefix identity at the unscaled query $\bar z$ gives
$\nabla f(z)=\nabla f^{[r+1]}(z)$.  Define
$\Delta^f_{r+1}(z)=\nabla f^{[r+1]}(z)-\nabla f^{[r]}(z)$.
The gradient multiplier in \eqref{eq:scaled-function} is $8\eps$, so
Lemma~\ref{lem:frontier-bound}, applied at $\bar z$, gives
$\norm{\Delta^f_{r+1}(z)}\le\Gamma$.  Consequently,
\[
 \E_B G(z;B)=\nabla f^{[r+1]}(z)=\nabla f(z),
\]
and
\[
 \E_B\norm{G(z;B)-\nabla f(z)}^2
 =\left(\frac1p-1\right)\norm{\Delta^f_{r+1}(z)}^2
 \le\left(\frac1p-1\right)\Gamma^2=\sigma^2.
\]
Freshness of $B$ gives the conditional statements.
\end{proof}

The construction uses exactly the oracle assumptions in
\eqref{eq:model-oracle}: a fixed unbiased vector-valued map with a uniform
second-moment bound.  It does not impose samplewise smoothness, or require
$G(\cdot;\xi)$ to be the gradient of a globally defined sample loss.  The
algorithm receives only the returned gradient vector and cannot reuse a
sample seed at a second query.

\section{Randomized lower bound}
\label{sec:haar}

\subsection{Concentration for adaptively queried frames}

We draw the primal frame $U$ and the $M$ dual frames $V_i$ independently from
their Haar distributions.  Let $d_b$ denote the ambient dimension of each
dual block, so that the total dual dimension is $d_y=Md_b$.  The proof uses
an ideal process whose revealed columns always form an initial segment of
\eqref{eq:global-order}.  Revealing a column means adjoining the entire
column to the proof's filtration; the algorithm itself still observes only
the prescribed oracle response.

We use the standard spherical tail bound: if $v$ is uniform on the unit
sphere of a $D$-dimensional subspace and $q$ is fixed, then
\begin{equation}
 \Pp\{|v^\top q|>a\}
 \le2\exp\!\left(-\frac{(D-1)a^2}{2\norm q^2}\right).
 \label{eq:spherical-tail}
\end{equation}
When $q=0$, the relevant probability is zero.

\begin{lemma}[Concentration for a prefix-revealing protocol]
\label{lem:adaptive-frame}
Consider a protocol that makes at most $Q$ queries and one final output
test.  Assume its next query is measurable with respect to its independent
random seeds and the columns already revealed in an initial segment of
\eqref{eq:global-order}.  Each new column is revealed only after that query
has been chosen, and the decision to reveal it does not inspect its value
or any later column.  If $0<\delta<1$ and
\begin{align}
 d_b&\ge2N+2^{20}N\log\frac{64MN(Q+1)}{\delta},
 \label{eq:dy-bound}\\
 d_x&\ge2M+2^{24}M\log\frac{64M(Q+1)}{\delta},
 \label{eq:dx-bound}
\end{align}
then, with probability at least $1-\delta$, at all these queries and at the
final output every still-unrevealed primal column satisfies
$|u_i^\top\varrho(\bar x)|\le1/4$, and every still-unrevealed dual column
satisfies \eqref{eq:weight-threshold} at $\bar z$ whenever its block is in
the hidden cone.
\end{lemma}

\begin{proof}
Condition on the protocol seeds used so far and on all previously revealed
columns.  By the stipulated revelation rule, the unused columns of each
frame remain Haar in the orthogonal complement of its revealed columns;
the unused parts of distinct frames remain conditionally independent.
The next query is fixed under this conditioning.

For a dual frame $V_i$, condition further on the entire primal frame and
all other dual frames.  This fixes
$\eta_i=\rho(u_{i-1}^\top\varrho(\bar x))$ without changing the conditional
Haar law of the unused part of $V_i$.  For $i=1$, the scale is the fixed
number $\rho(1)$.  If the block is in its hidden cone, then
$\norm{\bar y^{(i)}/\eta_i}\le32$, while
$c_*\sqrt{\omega_{j-1}}\ge1/(8\sqrt N)$.  The remaining dimension is at
least $d_b-N$.  Thus \eqref{eq:spherical-tail} bounds each dual failure
probability by
\[
 2\exp\!\left(-\frac{d_b-N-1}{2^{17}N}\right).
\]
If the block is public, no dual projection test is required.  The extra
conditioning is used only to estimate this probability; it is not added
to the protocol's filtration.

For the primal frame, use
$\norm{\varrho(\bar x)}\le230\sqrt M$, threshold $1/4$, and remaining
dimension at least $d_x-M$.  Each primal failure probability is at most
\[
 2\exp\!\left(-\frac{d_x-M-1}{1{,}692{,}800M}\right).
\]
The dimensions in \eqref{eq:dy-bound}--\eqref{eq:dx-bound} make the sum of
these bounds over at most $MN(Q+1)$ dual tests and $M(Q+1)$ primal tests
less than $\delta$.  The argument applies to each ideal query without
conditioning on earlier projection tests having passed.  A union bound
therefore proves the simultaneous statement, including the final output.
\end{proof}

\subsection{An ideal process and exact coupling}

\begin{lemma}[Exact prefix-based filtration coupling]
\label{lem:filtration}
Fix the algorithm's internal random seed.  There is an ideal process with
filtration $(\cF_t^{\rm id})$ and a revealed prefix of length $R_t$ such
that the following hold.
\begin{enumerate}
\item The next ideal query is $\cF_t^{\rm id}$-measurable.  Conditional on
this filtration, all unrevealed frame columns retain their Haar laws in
the appropriate orthogonal complements.
\item A Bernoulli failure reveals no new column.  A success can reveal
only the current inner dual column.  An exact terminal or connector
response can reveal only its current column.  A public-frontier response
reveals no column.
\item The ideal oracle makes branch decisions using only the query,
previously revealed columns, and a fresh Bernoulli seed; it never tests
an unrevealed projection.
\item On the simultaneous projection event in
Lemma~\ref{lem:adaptive-frame}, the ideal and true transcripts coincide
at every call, and hence have the same final output.
\end{enumerate}
\end{lemma}

\begin{proof}
Initialize the filtration with the public parameters and the algorithm's
fixed seed, and let $R_0=0$.  Suppose that positions
$1,\ldots,R_{t-1}$ have been revealed.  At the next ideal query, evaluate
all previously revealed latent projections at $\bar z=z/\lambda$,
replace all unrevealed latent projections by zero, and retain the public
radial norms.  Compute the resulting robust progress $r\le R_{t-1}$ and
public-block status.  Test the stronger projection thresholds only at
already revealed positions in $(r,R_{t-1}]$; provisionally treat every
unrevealed position as passing.

If $r=K$, all columns have been revealed and the ideal response is the
exact full gradient.  Otherwise, form the response as follows.
\begin{enumerate}
\item If a revealed later position fails its stronger threshold, return
$\nabla f^{[r+1]}(z)$.  Here $r+1\le R_{t-1}$, so no new column is used.
\item If the tests pass and $r+1$ is an inner dual position in a
nonpublic block, use \eqref{eq:fixed-oracle}.  On failure, the response
$\nabla f^{[r]}(z)$ uses only previously revealed columns.  On success,
adjoin the current column if $r+1=R_{t-1}+1$ and return the prescribed
prefix formula.  The decision to reveal precedes any inspection of
that column.
\item If the frontier is nonpublic and $r+1$ is a terminal dual
position, or if $r+1$ is a primal connector, adjoin
that column if it is new and return $\nabla f^{[r+1]}(z)$.
\item If the frontier is public, return $\nabla f^{[r]}(z)$ and reveal
no new column.
\end{enumerate}
Every prefix gradient is, by definition, a function only of the columns
in that prefix, even when a projection test would fail.  Thus each ideal
response is measurable with respect to the old filtration and at most
one newly revealed column.  Append the response, the fresh seed, and any
new column to the filtration.  The decision to reveal a new column uses
only the old filtration and the fresh seed, so the conditional Haar law
is preserved.  This proves that the ideal process satisfies the
hypotheses of Lemma~\ref{lem:adaptive-frame}.

For transcript equality, work on that lemma's simultaneous event.
Every unrevealed coordinate lies in its robust dead zone and passes its
stronger threshold.  Replacing these coordinates by zero leaves the
robust progress unchanged.  It also leaves public-block status
unchanged: if a predecessor primal column is unrevealed, then
$|s_{i-1}|\le1/4$ and $\rho(s_{i-1})=0$; otherwise its projection is
retained exactly.  The true simultaneous test therefore fails exactly
when the ideal test detects a failing revealed position.

If the test fails, all positions after $r$ are nevertheless in their
robust dead zones, by the definition of $r$.  The prefix identity gives
$\nabla f(z)=\nabla f^{[r+1]}(z)$, so the two exact responses agree.
On a passing inner frontier they use identical prefixes and the same
Bernoulli seed.  At terminal and connector frontiers the boundary
identities give the same equality with the exact full gradient.  At a
public frontier the current block has no hidden dependence, and all
later blocks have zero predecessor scale; hence the exact gradient
equals $\nabla f^{[r]}(z)$.  The $r=K$ case is immediate.
Induction proves equality at every call.  The final output is a
measurable function of the common transcript and is covered by the
additional projection test.
\end{proof}

\begin{lemma}[One-coordinate progress for unrestricted algorithms]
\label{lem:progress}
In the ideal process, each call reveals at most one new position in
\eqref{eq:global-order}.  Each inner dual position $(i,j)$ with $j<N$
requires a Bernoulli success before it can be revealed.  Consequently,
revealing the terminal primal column $u_M$ requires at least
\begin{equation}
 S=M(N-1)
 \label{eq:S}
\end{equation}
Bernoulli successes.  On the simultaneous projection event, the same
information restriction applies to the true transcript and its final
output.
\end{lemma}

\begin{proof}
The revealed set is an initial segment.  At a new inner frontier a
failed stronger test is impossible in the ideal process, since any
position that could fail that test is already revealed.  A public
frontier reveals no column, and a Bernoulli failure returns the previous
prefix gradient.  Thus only a success reveals a new inner column.
Exact terminal and connector responses add at most their own column
and cannot bypass an inner position.  The initial segment ending at
$u_M$ contains precisely $M(N-1)$ inner positions.  Transcript equality
and the output projection test give the final assertion.
\end{proof}
\subsection{Completion of the proof of Theorem~\ref{thm:main-lower}}
\label{sec:completion}

Fix an arbitrary algorithm, including its measurable query and output
maps in every finite dimension. Let $S=M(N-1)$ and choose
\begin{equation}
 Q=\left\lfloor\frac{S}{16p}\right\rfloor,
 \qquad
 p=\frac{\Gamma^2}{\Gamma^2+\sigma^2}.
 \label{eq:Qchoice}
\end{equation}
Because $S/p\ge S\ge56$, the floor satisfies
\begin{equation}
 Q\ge\frac{S}{32p}.
 \label{eq:Q-floor-lower}
\end{equation}
Any algorithm using fewer calls may be padded by query repetitions, so it is enough to prove failure through this horizon.  Choose $d_b$ and $d_x$ according to
\eqref{eq:dy-bound}--\eqref{eq:dx-bound} with $\delta=1/16$, and let
$\mathcal E$ be the simultaneous projection event for the ideal
process, including its final output.  Draw independent Bernoulli
variables $B_1,\ldots,B_Q$ at every call, including the exact calls on
which they are unused.  The total number of newly revealed inner
columns is at most $\sum_{t=1}^Q B_t$.  Lemma~\ref{lem:progress}
therefore implies
\[
 \mathcal E\cap\{\text{the terminal primal column is revealed}\}
 \subseteq
 \left\{\sum_{t=1}^Q B_t\ge S\right\}.
\]
The Bernoulli sum is binomial unconditionally; no independence after
conditioning on $\mathcal E$ is needed.  Markov's inequality yields
\begin{equation}
 \Pp\!\left\{\sum_{t=1}^Q B_t\ge S\right\}
 \le\frac{pQ}{S}\le\frac1{16}.
 \label{eq:markov-progress}
\end{equation}
With probability at least $1-1/16-1/16=7/8$, the event
$\mathcal E$ holds and the terminal primal column is still unrevealed.
The final-output test then gives
$|u_M^\top\varrho(\widehat x/\lambda)|\le1/4$, hence
$\prog_{1/4}(U^\top\varrho(\widehat x/\lambda))<M$.
Transcript equality transfers this conclusion to the true oracle.
Lemma~\ref{lem:scaled-obstruction} now gives
\[
 \E\!\left[
  \norm{\nabla_xf(\widehat x,\widehat y)}^2
  +\kappa\norm{\nabla_yf(\widehat x,\widehat y)}^2
 \right]\ge\frac78\,4\eps^2>\eps^2.
\]

Finally,
\begin{equation}
 \frac1p=1+\frac{\sigma^2}{\Gamma^2}
 =1+\frac{N\sigma^2}{64^2\eps^2},
 \label{eq:pinv}
\end{equation}
so, since $M,N\ge8$,
\begin{align}
 \frac{S}{p}
 &\ge c\left(MN+\frac{MN^2\sigma^2}{\eps^2}\right)\notag\\
 &\ge c'\left(
  \frac{L\kappa H}{\eps^2}
  +\frac{L\kappa^2H\sigma^2}{\eps^4}
 \right),
 \label{eq:final-rate}
\end{align}
where $c,c'>0$ are universal because $c_M,c_N$ are fixed numerical constants.

The construction parameters and frame distribution are independent of
the algorithm's realized seed. For every fixed seed, the same coupling
and bounds apply; averaging over that seed therefore preserves the
displayed expectation bound, jointly over the Haar frames, the
algorithm's internal seed, and the fresh oracle seeds. Hence some fixed
realization of all hidden frames has expected residual greater than
$\eps^2$, where only the algorithm and oracle remain random. This
realization determines a deterministic hard objective and a fixed
oracle map. Together with \eqref{eq:Q-floor-lower} and
\eqref{eq:final-rate}, this proves Theorem~\ref{thm:main-lower}.

\section{Discussion}
\label{sec:discussion}

The matching rate in \eqref{eq:matching-rate} identifies the optimal
dependence on the condition parameter for weighted joint stationarity
in the fresh unbiased-gradient model. The deterministic term reflects
the number of coordinates in the interleaved saddle chain. The
stochastic term also reflects the size of the full gradient signal
that reveals one coordinate. Controlling that full signal, rather than
one selected gradient component, is what makes the variance calculation
compatible with the perspective construction.

The information argument is separated from the analytical properties
of the hard family. The radial cap preserves global smoothness and
the dual P\L{} inequality while bounding the region that can expose a
hidden frame. Prefix functions define a measurable ideal transcript
for every query, and concentration is used only to identify that
transcript with the fixed true oracle. This order of definitions is
essential: a response that is independent of future directions only
on a favorable event does not by itself justify a deferred-decision
argument.

The theorem has three specific scope boundaries. Its stationarity
target is the weighted residual of $f$; its initialization parameter
is the charged budget $H$; and its oracle is a fixed unbiased vector
oracle with independent fresh calls. Matching rates for value-gradient
stationarity with an uncharged value gap, for prescribed small
dimensions, or for additional sample-function regularity require
separate arguments. The present dimension bounds and numerical
constants establish existence and uniformity; they are not optimized.

\clearpage
\appendix
\section{Derivative bounds for the outer scale}
\label{app:rho}

For $s>1/2$, put $t=2s-1>0$.  Then
\[
 \rho(s)=\exp\!\left(\frac12-\frac1{2t^2}\right),
 \quad
 \rho'(s)=2\rho(s)t^{-3},
 \quad
 \rho''(s)=4\rho(s)t^{-6}-12\rho(s)t^{-4}.
\]
Writing $u=t^{-2}$ gives
$|\rho'|=2e^{1/2}u^{3/2}e^{-u/2}$, whose maximum occurs at $u=3$ and is below $4$.  Likewise,
$|\rho''|=4e^{1/2}e^{-u/2}u^2|u-3|$, whose maximum is below $40$.  The negative branch is symmetric and all derivatives vanish at the flat-zone interfaces.  This proves \eqref{eq:rho-bounds}.

\section{Explicit constants and dimensions}
\label{app:constants}

The table collects constants used by the construction. All are
independent of $M,N,d_x,d_b$ and the target problem parameters.
The per-block dimension is $d_b$ and the total dual dimension is $Md_b$.


\begin{longtable}{>{\raggedright\arraybackslash}p{0.33\textwidth}p{0.57\textwidth}}
\toprule
Quantity & Value or bound \\
\midrule
\endfirsthead
\toprule
Quantity & Value or bound \\
\midrule
\endhead
Shift and cap parameters & $\tau=1/4$, $R=16$, $A=1$, $C=32$, $a\in[0,10]$ \\
Scalar derivatives & $\norm{q_\tau'}_\infty\le2$, $\Lip(q_\tau')\le32/3$, $\norm{p'}_\infty\le3/4$, $\Lip(p')\le3/2$ \\
Dual-chain P\L{} and smoothness & $\norm{\nabla D_a}^2\ge D_a/(385N)$; $\Lip(\nabla D_a)\le128$ \\
Radius-$32$ robust-chain bounds & $D_a\le1040$, $\norm{\nabla D_a}\le4097$ \\
Cap P\L{} and smoothness & $\norm{\nabla\mathscr D}^2\ge10^{-5}\mathscr D/N$; $\Lip(\nabla\mathscr D)\le1024$ \\
Local block & $L_C=4\cdot10^7$; primal-gradient norm $\le2\cdot10^6$ \\
Full unscaled saddle & $\ell_0=10^8$; dual P\L{} modulus $\mu_0/N$ with $\mu_0=5\cdot10^{-6}$ \\
Transfer coefficient & $L_{\rm def}=5\cdot10^7$, $K_{\rm tr}^2=2\cdot10^{13}$ \\
Chain lengths & $N=\lfloor c_N\kappa\rfloor$, $M=\lfloor c_MLH/\eps^2\rfloor$ with \eqref{eq:cN-cM} \\
Global scaling & $\lambda=8\ell_0\eps/L$; gradient multiplier $8\eps$ \\
Inner stochastic signal & unscaled $\le8/\sqrt N$; scaled $\Gamma=64\eps/\sqrt N$ \\
Frame dimensions & \eqref{eq:dy-bound} and \eqref{eq:dx-bound} \\
Failure probabilities & projection-test failure $\le1/16$; enough Bernoulli successes $\le1/16$ \\
\bottomrule
\end{longtable}

\begin{thebibliography}{9}

\bibitem{arjevani}
Yossi Arjevani, Yair Carmon, John C. Duchi, Dylan J. Foster,
Nathan Srebro, and Blake Woodworth.
Lower bounds for non-convex stochastic optimization.
\emph{Mathematical Programming}, 199:165--214, 2023.
\href{https://doi.org/10.1007/s10107-022-01822-7}{doi:10.1007/s10107-022-01822-7}.

\bibitem{carmon}
Yair Carmon, John C. Duchi, Oliver Hinder, and Aaron Sidford.
Lower bounds for finding stationary points I.
\emph{Mathematical Programming}, 184:71--120, 2020.
\href{https://doi.org/10.1007/s10107-019-01406-y}{doi:10.1007/s10107-019-01406-y}.

\bibitem{panli}
Siyu Pan and Jiajin Li.
Lower bounds for nonconvex--P\L{} minimax optimization.
\emph{arXiv preprint arXiv:2608.26799}, 2026.
\url{https://arxiv.org/abs/2608.26799}.

\bibitem{yang}
Junchi Yang, Antonio Orvieto, Aurelien Lucchi, and Niao He.
Faster single-loop algorithms for minimax optimization without strong concavity.
In \emph{Proceedings of the 25th International Conference on Artificial
Intelligence and Statistics}, volume 151 of \emph{Proceedings of Machine
Learning Research}, pages 5485--5517, 2022.
\url{https://proceedings.mlr.press/v151/yang22b.html}.

\end{thebibliography}
\end{document}
